\section{Contoh Perhitungan Hash (Simulasi)}

\subsection{Konsep Fungsi Kompresi}

Fungsi hash modern (seperti SHA-256) bekerja dengan memproses pesan dalam blok-blok menggunakan \textbf{fungsi kompresi}. Untuk memahami ini, mari kita lakukan simulasi manual satu langkah fungsi kompresi sederhana.

\textbf{Algoritma Mainan (Toy Hash)}:
\begin{itemize}
  \item \textbf{Input}: Pesan 4-bit ($M$) dan State 4-bit ($H$)
  \item \textbf{Output}: State baru 4-bit ($H'$)
  \item \textbf{Fungsi Kompresi}: $H' = (H \oplus M) + (H \lll 1) \pmod{16}$
\end{itemize}

\subsection{Simulasi Step-by-Step}

\textbf{Pesan}: "A" (ASCII 65 = 0100 0001 biner)
\textbf{IV (Initial Vector)}: 0000

Pesan dibagi menjadi 2 blok 4-bit:
\begin{itemize}
  \item $M_1 = 0100$ (4)
  \item $M_2 = 0001$ (1)
\end{itemize}

\textbf{Langkah 1: Proses Blok 1 ($M_1 = 4$)}
\begin{itemize}
  \item $H_{in} = 0000$ (0)
  \item $H \oplus M = 0000 \oplus 0100 = 0100$ (4)
  \item $H \lll 1 = 0000 \text{ rotasi kiri 1} = 0000$ (0)
  \item $H_{out} = (4 + 0) \pmod{16} = 4$ (0100)
\end{itemize}
State sekarang: $H_1 = 0100$

\textbf{Langkah 2: Proses Blok 2 ($M_2 = 1$)}
\begin{itemize}
  \item $H_{in} = 0100$ (4)
  \item $H \oplus M = 0100 \oplus 0001 = 0101$ (5)
  \item $H \lll 1 = 0100 \text{ rotasi kiri 1} = 1000$ (8)
  \item $H_{out} = (5 + 8) \pmod{16} = 13$ (1101)
\end{itemize}
State sekarang: $H_2 = 1101$ (D dalam heksadesimal)

\textbf{Hash Akhir}: D

\subsection{Efek Avalanche}

Mari ubah 1 bit pesan awal. "A" (0100 0001) $\rightarrow$ "C" (0100 0011).
Perubahan pada blok 2: $M'_2 = 0011$ (3).

\textbf{Langkah 1 (Sama)}: $H_1 = 0100$

\textbf{Langkah 2 (Baru)}:
\begin{itemize}
  \item $H_{in} = 0100$ (4)
  \item $H \oplus M' = 0100 \oplus 0011 = 0111$ (7)
  \item $H \lll 1 = 1000$ (8)
  \item $H_{out} = (7 + 8) \pmod{16} = 15$ (1111)
\end{itemize}
State Akhir: F

Perubahan 1 bit input mengubah output dari D (1101) menjadi F (1111). Meskipun sederhana, ini mengilustrasikan efek avalanche.

\subsection{Struktur Logika SHA-256 Satu Langkah}

SHA-256 jauh lebih kompleks. Satu langkah (step) dari 64 langkah melibatkan:

\[
T_1 = h + \Sigma_1(e) + Ch(e,f,g) + K_t + W_t
\]
\[
T_2 = \Sigma_0(a) + Maj(a,b,c)
\]
\[
h = g, g = f, f = e, e = d + T_1, d = c, c = b, b = a, a = T_1 + T_2
\]

Dimana:
\begin{itemize}
  \item $\Sigma_0, \Sigma_1$: Rotasi dan shift bit (diffusion)
  \item $Ch, Maj$: Fungsi logik non-linear (confusion)
  \item $W_t$: Message schedule (perluasan input)
  \item $K_t$: Konstanta ronde (untuk menghambat analisis)
\end{itemize}

Setiap operasi di atas adalah operasi 32-bit. Kompleksitas ini memastikan bahwa tidak ada jalan pintas matematis untuk membalikkan fungsi hash.
